\section{Section a) Ground-breaking Nature \& Objectives}

\subsection{State of the Art and Limitations}
Current paradigms in quantum error correction (QEC) and topological quantum computing assume that the physical environment is memoryless (Markovian). Under this assumption, passive topological protection (e.g., in toric codes or Majorana zero modes) offers an elegant path to fault tolerance. In practice, however, superconducting, spin-orbit, and trapped-ion systems are coupled to environments featuring $1/f$ charge or flux noise, nuclear spin baths, and phonon-mediated losses. These sources induce strong non-Markovian memory effects, creating temporal correlations that break the standard code assumptions. 

\pbreak
Consequently, current topological states are highly vulnerable to high-frequency spectral components and low-frequency drift. There is currently no mathematical framework that describes topological phase transitions under non-Markovian, time-dependent stochastic fields, nor exists an experimental methodology to actively use environmental noise to stabilize topological structures.

\subsection{Project Objectives}
QUPROB proposes to transform our understanding of open topological quantum systems. Instead of treating environmental noise solely as an adversary, we propose to exploit and engineer the stochastic environment to actively preserve topological invariants. The project is organized around four major objectives:

\pbreak
\textbf{Objective 1: Establish Stochastic Topological Field Theory (STFT).} 
Develop a rigorous mathematical formulation of topological invariants (such as the Chern number and Wannier centers) for non-equilibrium systems governed by non-Markovian stochastic differential equations.
\pbreak
\textbf{Objective 2: Discover Noise-Induced Topological Phase Transitions.}
Identify and classify phase boundaries where stochastic fluctuations dynamically induce, rather than destroy, topological order. We will exploit the phenomenon of quantum stochastic resonance to lock edge states in place.
\pbreak
\textbf{Objective 3: Develop Active Stochastic Control Protocols.}
Design high-frequency Floquet engineering and dynamic pulse sequences that shift environmental spectral densities away from critical system transitions, creating an active shield.
\pbreak
\textbf{Objective 4: Experimental Proof-of-Concept.}
Validate these protocols on a superconducting qubit processor at TUM, demonstrating a stochastically protected logical qubit with a lifetime exceeding the physical limit of its constituents.

\vspace{1.5em}
